\ticket{3}{Основные средства современных функциональных языков (на примере языка Хаскель)}

\begin{examplan}
    \item Кратко охарактеризовать Haskell как чисто функциональный язык.
    \item Объяснить чистые функции, ссылочную прозрачность и управление эффектами.
    \item Рассказать о функциях высшего порядка, каррировании и частичном применении.
    \item Описать статическую типизацию, вывод типов, алгебраические типы данных и классы типов.
    \item Объяснить ленивые вычисления, сопоставление с образцом и монады.
\end{examplan}

\subsection*{Короткая версия для письменного ответа}

\textbf{Haskell} --- чисто функциональный язык программирования общего назначения со статической типизацией, выводом типов и ленивыми вычислениями. В Haskell программа строится как набор выражений и функций, а не как последовательность команд, изменяющих состояние.

Основные средства Haskell:
\begin{itemize}
    \item чистые функции и ссылочная прозрачность;
    \item функции высшего порядка;
    \item каррирование и частичное применение;
    \item статическая типизация и вывод типов;
    \item параметрический полиморфизм;
    \item ленивые вычисления;
    \item сопоставление с образцом;
    \item алгебраические типы данных;
    \item классы типов;
    \item монады для описания вычислений с контекстом и эффектами.
\end{itemize}

\subsection*{Чистота и ссылочная прозрачность}

\textbf{Чистая функция} при одинаковых аргументах всегда возвращает одинаковый результат и не имеет побочных эффектов: не изменяет глобальные переменные, не выполняет ввод-вывод напрямую, не меняет внешнее состояние.

\textbf{Ссылочная прозрачность} означает, что выражение можно заменить его значением без изменения поведения программы. Например, если \texttt{square 3} равно \texttt{9}, то в чистой программе это выражение можно заменить на \texttt{9}.

\textbf{Преимущества чистоты:}
\begin{itemize}
    \item легче рассуждать о корректности программы;
    \item функции проще тестировать изолированно;
    \item компилятору проще выполнять оптимизации;
    \item отсутствие общего изменяемого состояния облегчает параллелизм.
\end{itemize}

Побочные эффекты в Haskell не исчезают, а явно отражаются в типах. Например, \texttt{IO String} означает действие ввода-вывода, которое при выполнении возвращает строку. Благодаря этому чистая часть программы отделяется от взаимодействия с внешним миром.

\subsection*{Функции высшего порядка}

В Haskell функции являются объектами первого класса: их можно передавать как аргументы, возвращать из других функций и связывать с именами.

\textbf{Функция высшего порядка} --- функция, которая принимает другую функцию как аргумент или возвращает функцию как результат.

\begin{verbatim}
map    :: (a -> b) -> [a] -> [b]
filter :: (a -> Bool) -> [a] -> [a]

map (+1) [1,2,3]        -- [2,3,4]
filter even [1,2,3,4]  -- [2,4]
\end{verbatim}

Композиция функций позволяет собирать сложное преобразование из простых:

\begin{verbatim}
(not . even) 3  -- True
\end{verbatim}

Такие средства делают код декларативным: программа описывает цепочку преобразований данных.

\subsection*{Каррирование и частичное применение}

\textbf{Каррирование} --- представление функции от нескольких аргументов как последовательности функций от одного аргумента. В Haskell все функции формально принимают один аргумент.

\begin{verbatim}
add :: Int -> Int -> Int
add x y = x + y
\end{verbatim}

Тип \texttt{Int -> Int -> Int} читается как \texttt{Int -> (Int -> Int)}: функция получает первый \texttt{Int} и возвращает функцию, ожидающую второй \texttt{Int}. Поэтому \texttt{add 5 3} интерпретируется как \texttt{(add 5) 3}.

\textbf{Частичное применение} --- создание новой функции путем фиксации части аргументов:

\begin{verbatim}
addFive :: Int -> Int
addFive = add 5

addFive 3  -- 8
\end{verbatim}

\subsection*{Статическая типизация и вывод типов}

Haskell имеет строгую статическую типизацию: типы выражений проверяются на этапе компиляции. Многие ошибки обнаруживаются до запуска программы.

Система типов Haskell основана на идеях Хиндли--Милнера и поддерживает \textbf{вывод типов}: компилятор часто сам выводит наиболее общий тип выражения.

\begin{verbatim}
identity :: a -> a
identity x = x

apply :: (a -> b) -> a -> b
apply f x = f x
\end{verbatim}

Буквы \texttt{a} и \texttt{b} обозначают переменные типов. Это пример \textbf{параметрического полиморфизма}: функция работает для любых типов, если ее определение не требует конкретных операций над ними.

Явные аннотации типов часто не обязательны, но считаются хорошим стилем: они документируют намерение программиста и помогают найти ошибки.

\subsection*{Ленивые вычисления}

\textbf{Ленивые вычисления} означают, что выражение вычисляется только тогда, когда его результат действительно нужен. Такой способ называют \textit{call-by-need}.

\textbf{Преимущества ленивости:}
\begin{itemize}
    \item можно работать с бесконечными структурами данных;
    \item промежуточные структуры не всегда вычисляются полностью;
    \item управляющие конструкции можно задавать как обычные функции.
\end{itemize}

\begin{verbatim}
ones :: [Int]
ones = 1 : ones

take 5 ones  -- [1,1,1,1,1]
\end{verbatim}

\textbf{Недостатки ленивости:} иногда сложнее предсказать порядок вычислений, использование памяти и производительность.

\subsection*{Сопоставление с образцом}

\textbf{Сопоставление с образцом} позволяет анализировать структуру данных и извлекать ее части прямо в определениях функций.

\begin{verbatim}
factorial :: Int -> Int
factorial 0 = 1
factorial n = n * factorial (n - 1)

length' :: [a] -> Int
length' []     = 0
length' (_:xs) = 1 + length' xs
\end{verbatim}

Сопоставление с образцом делает код компактным и наглядным: разные формы входных данных описываются разными уравнениями функции.

\subsection*{Алгебраические типы данных}

\textbf{Алгебраические типы данных} позволяют строить собственные типы из вариантов и наборов полей.

\textbf{Тип-сумма} задает выбор между несколькими конструкторами:

\begin{verbatim}
data Shape
    = Circle Float
    | Rectangle Float Float

area :: Shape -> Float
area (Circle r)      = pi * r * r
area (Rectangle w h) = w * h
\end{verbatim}

\textbf{Тип-произведение} объединяет несколько значений в одну структуру:

\begin{verbatim}
data Point = Point Float Float

data Person = Person
    { name :: String
    , age  :: Int
    }
\end{verbatim}

Важные стандартные ADT:

\begin{verbatim}
data Maybe a = Nothing | Just a
\end{verbatim}

\texttt{Maybe a} используется, когда результат может отсутствовать: \texttt{Just x} означает наличие значения, \texttt{Nothing} --- отсутствие значения.

\subsection*{Классы типов}

\textbf{Класс типов} в Haskell --- механизм ad-hoc полиморфизма. Он задает набор операций, которые должен поддерживать тип, чтобы быть экземпляром этого класса.

Классы типов похожи на интерфейсы по идее контракта, но они описывают поведение типов, а не объектов.

\textbf{Примеры стандартных классов типов:}
\begin{itemize}
    \item \texttt{Eq}: сравнение на равенство, операции \texttt{==} и \texttt{/=};
    \item \texttt{Ord}: упорядочивание;
    \item \texttt{Show}: преобразование значения в строку;
    \item \texttt{Read}: чтение значения из строки;
    \item \texttt{Num}: числовые операции;
    \item \texttt{Functor}, \texttt{Applicative}, \texttt{Monad}: работа с вычислениями в контексте.
\end{itemize}

Ограничения на типы записываются через \texttt{=>}:

\begin{verbatim}
areEqual :: Eq a => a -> a -> Bool
areEqual x y = x == y
\end{verbatim}

Эта функция работает для любого типа \texttt{a}, если для него определено сравнение на равенство.

\subsection*{Монады}

\textbf{Монада} --- это класс типов, который задает способ последовательно связывать вычисления с некоторым контекстом: возможной ошибкой, отсутствием значения, списком вариантов, состоянием или вводом-выводом.

Главная идея монады: результат одного шага передается в следующий шаг, при этом контекст вычисления сохраняется и обрабатывается единообразно.

Основные операции:
\begin{itemize}
    \item \texttt{return} или \texttt{pure}: помещает обычное значение в контекст;
    \item \texttt{>>=}: связывает монадическое значение со следующей функцией.
\end{itemize}

\textbf{Примеры монад:}
\begin{itemize}
    \item \texttt{Maybe}: вычисление может не дать результата;
    \item \texttt{Either e}: вычисление может завершиться ошибкой;
    \item список \texttt{[]}: вычисление может иметь несколько результатов;
    \item \texttt{IO}: действия ввода-вывода;
    \item \texttt{State}: вычисления с состоянием.
\end{itemize}

\begin{verbatim}
main :: IO ()
main = do
    putStrLn "What is your name?"
    name <- getLine
    putStrLn ("Hello, " ++ name ++ "!")
\end{verbatim}

\texttt{do}-нотация является удобной записью последовательности монадических действий. В случае \texttt{IO} она позволяет описывать ввод-вывод, не нарушая чистоту обычных функций: эффект явно виден в типе.

\begin{important}
    \item Haskell --- чисто функциональный язык: вычисление строится вокруг выражений и функций.
    \item Чистота означает отсутствие неявных побочных эффектов; эффекты отражаются в типах, например \texttt{IO a}.
    \item Функции высшего порядка, каррирование и частичное применение дают мощные средства абстракции.
    \item Статическая типизация и вывод типов позволяют писать общий, но типобезопасный код.
    \item ADT и pattern matching являются основными средствами моделирования и разбора данных.
    \item Type classes реализуют ad-hoc полиморфизм, а монады позволяют единообразно описывать вычисления с контекстом.
\end{important}
